\documentclass[11pt]{article}
\usepackage[margin=1in]{geometry}
\usepackage{hyperref}
\usepackage{amsmath, amsthm, amssymb}
\usepackage{xcolor}
\usepackage{tcolorbox}
\tcbset{colback=gray!5, colframe=gray!50, boxrule=0.4pt, sharp corners,
        before skip=4pt, after skip=4pt}

\title{Citation Verification for \\
``Surrogate Inverse-Variance Weighting for Intervention Harvesting
on Heavy-Tailed Click Logs''}
\author{(internal verification document, not submitted)}
\date{}

\newcommand{\paperclaim}[1]{\textbf{Paper claim:} #1}
\newcommand{\source}[1]{\textbf{Source location:} #1}
\newcommand{\quoted}[1]{\begin{tcolorbox}\textit{#1}\end{tcolorbox}}

\begin{document}
\maketitle

\noindent This document quotes verbatim from each cited paper the passage
that supports the corresponding claim in the CIKM short paper. Goal: let
the authors verify that each citation accurately backs the claim it is
attached to. Local PDFs of the cited papers live in
\texttt{paper/cikm/citations/}.

\section*{Status summary}

\noindent\textbf{Fully verified} (author, title, and load-bearing claim
confirmed by quoting the relevant passage from the local PDF):
\texttt{agarwal2019estimating}, \texttt{borenstein2009meta},
\texttt{bottou2013counterfactual} (§4.4 heavy-tailed-weights variance),
\texttt{chuklin2015clickmodels}, \texttt{craswell2008experimental}
(cascade model definition), \texttt{hesterberg1995},
\texttt{joachims2017unbiased}, \texttt{kong1992ess} (via
Owen~Eq.~9.13), \texttt{owen2013mc}, \texttt{swaminathan2015snips}
(correct PDF in \texttt{swaminathan\_joachims\_2015\_nips\_snips.pdf}),
\texttt{wang2018position} (verified for EM-style estimation; see action
item below), \texttt{yue2010beyond} (correct PDF in
\texttt{yue\_2010\_www\_beyond.pdf}), \texttt{zou2022large} (verified
size and \emph{public-aggregation} framing; the $94\%/78\%$ statistics
are our own computation).

\noindent\textbf{Author/title only verified} (partial PDF, content not in
local copy): \texttt{vdv1998asymptotic} (frontmatter only; equivalent
delta-method material in Wasserman 36-705 lec~3 which we have).

\noindent\textbf{No PDF expected} (dataset citation): \texttt{kddcup2012}
(KDD Cup 2012 Track 2 competition release).

\medskip

\noindent\textbf{Action taken in the paper based on this verification:}
\begin{itemize}
\item Removed \texttt{wang2018position} from the ``PBM violations on real
logs are documented in~[\,...]'' cite list, since Wang 2018 works
\emph{within} PBM (its EM algorithm estimates PBM parameters; it does
not empirically demonstrate PBM violation). Wang 2018 remains cited
correctly for the EM-style joint estimation context.
\item Dropped \texttt{aitken1934least} and \texttt{hedges1985meta} (no
freely downloadable PDF) in favour of \texttt{borenstein2009meta} which
covers the same IVW material.
\end{itemize}

\bigskip

\section*{1. \texttt{agarwal2019estimating}}
\paperclaim{The intervention-harvesting estimator of Agarwal et al.\ 2019
uses uniform per-cell weighting (each cell's click rate enters the sum with
weight 1).}\\
\source{Section 4.1 ``Local Estimators'', PivotOne and AdjacentChain
definitions on p.~4--5 of the WSDM 2019 paper.}

\quoted{Since the click counts $\hat c_k^{k,k'}$ and $\hat c_{k'}^{k,k'}$
can be treated just like data from an explicit intervention, the same
estimators apply. In particular, we observe the interventional set
$S_{1,k}$ and can thus use the same estimator that we previously used
for the explicit Swap(1,$k$) experiment. We call this the \textbf{PivotOne}
estimator
\[\hat p_k/\hat p_1 = \hat c_k^{1,k} / \hat c_1^{1,k}\,.\]
We can similarly adapt the estimator used in [29], which uses a chain of
swaps between adjacent positions in the ranking. We call this the
\textbf{AdjacentChain} estimator
$\hat p_k/\hat p_1 = (\hat c_2^{1,2}/\hat c_1^{1,2}) \cdot
(\hat c_3^{2,3}/\hat c_2^{2,3}) \cdots
(\hat c_k^{k-1,k}/\hat c_{k-1}^{k-1,k})$.}

\noindent Each $\hat c$ is defined earlier (Section 3.4) as
\[\hat c_k^{k,k'} := \sum_{i=1}^m \sum_{j=1}^{n_i} \sum_{d\in\Omega(q_i^j)}
\mathbf 1[(q_i^j,d)\in S_{k,k'}]\,\mathbf 1[rk(d|y_i^j)=k]\;
\frac{c_i^j(d)}{w(q_i^j,d,k)}.\]
The weight $w(q,d,k)$ normalises by ranker-assignment count
(propensity-style), and at the cell aggregation level the per-cell click
rate $C_k^i/N_k^i$ enters with coefficient~1 (no heteroscedastic re-weighting).
\textbf{No alternative per-cell weighting (min, harmonic, ...) is discussed
in the paper.} Verified.

\bigskip

\section*{2. \texttt{borenstein2009meta}}
\paperclaim{(i) Equal-weighting heteroscedastic estimators is statistically
suboptimal by classical inverse-variance theory; (ii) inverse-variance
weighting is canonical in meta-analysis; (iii) the count-only IVW optimum
restricted to count-only weights is the harmonic mean (a classical result
that Borenstein presents in standard form).}\\
\source{Borenstein, Hedges, Higgins, Rothstein, \emph{Introduction to
Meta-Analysis}, Wiley 2009. Sample chapter PDF in
\texttt{citations/borenstein\_2009\_meta\_chapter.pdf}.}

\quoted{Borenstein et al.\ Ch.~11 summary (p.~67): ``Under the fixed-effect
model all studies in the analysis share a common true effect. The summary
effect is our estimate of this common effect size, and the null hypothesis
is that this common effect is zero (for a difference) or one (for a
ratio). All observed dispersion reflects sampling error, and \textbf{study
weights are assigned with the goal of minimizing this within-study error.}''}

\quoted{Ch.~12 (p.~72), under ``Performing a random-effects meta-analysis''
(the same logic applies to fixed-effect meta-analysis with within-study
variance only): ``In order to obtain the most precise estimate of the
overall mean (to minimize the variance) we compute a weighted mean,
\textbf{where the weight assigned to each study is the inverse of that
study's variance.}''}

\quoted{Ch.~13 (p.~78, ``Estimating the Summary Effect''): ``Under the
fixed-effect model we assume that the true effect size for all studies is
identical, and the only reason the effect size varies between studies is
sampling error (error in estimating the effect size). Therefore, when
assigning weights to the different studies we can largely ignore the
information in the smaller studies since we have better information about
the same effect size in the larger studies.''}

\noindent Verified. The cite supports both ``equal-weighting heteroscedastic
estimators is statistically suboptimal'' and ``inverse-variance weighting is
canonical in meta-analysis''. (Aitken 1934 is the original GLS / BLUE
source; Borenstein covers the same content in modern textbook form.)

\bigskip

\section*{3. \texttt{bottou2013counterfactual}}
\paperclaim{The closest off-policy analogue to our plug-in $\widetilde V$
framework is SNIPS-style variance analysis of self-normalised IPS in
counterfactual evaluation; weight heterogeneity (not sample count) governs
effective precision on logged feedback.}\\
\source{Bottou, Peters, Quiñonero-Candela, Charles, Chickering, Portugaly,
Ray, Simard, Snelson, ``Counterfactual Reasoning and Learning Systems:
The Example of Computational Advertising'', JMLR 14:3207--3260, 2013.
PDF: \texttt{citations/bottou\_2013\_jmlr.pdf}. Author/title verified
from p.~1.}

\quoted{From the abstract: ``This work shows how to leverage causal
inference to understand the behavior of complex learning systems
interacting with their environment and predict the consequences of
changes to the system. [...] illustrated by experiments on the ad
placement system associated with the Bing search engine.''}

\quoted{\S4.4 Confidence Intervals (p.~3225): ``Unfortunately, when
$P(\omega)$ is small, the reweighting ratio $w(\omega)$ takes large
values with low probability. \textbf{This heavy tailed distribution has
annoying consequences because the variance of the integrand could be
very high or infinite.} [...] Importance sampling works best when the
actual distribution and the counterfactual distribution overlap. When
the counterfactual distribution has significant mass in domains where
the actual distribution is small, \textbf{the few samples available in
these domains receive very high weights. Their noisy contribution
dominates the reweighted estimate.}''}

\noindent Verified. \S4.4 (``Confidence Intervals'') explicitly analyses
how heavy-tailed importance weights dominate the variance of the
self-normalised IPS estimator -- exactly the off-policy analogue our
paper invokes.

\bigskip

\section*{4. \texttt{chuklin2015clickmodels}}
\paperclaim{Click-model survey including PBM and other click models;
canonical reference for click-model background.}\\
\source{Chuklin, Markov, de Rijke, \emph{Click Models for Web Search},
Morgan \& Claypool 2015 (author's version PDF in
\texttt{citations/chuklin\_2015\_clickmodels.pdf}).
Table of contents confirms PBM and Cascade among the basic click models.}

\quoted{From the Contents (p.~ii): ``Chapter~3 Basic Click Models:
\textbf{3.3 Position-based Model (PBM)}, \textbf{3.4 Cascade Model (CM)},
3.5 User Browsing Model (UBM), 3.6 Dependent Click Model (DCM), 3.7 Click
Chain Model (CCM), 3.8 Dynamic Bayesian Network Model (DBN).''}

\noindent Verified. The book is the canonical click-model survey and
explicitly covers PBM and its alternatives. Use to back the claim
``click models beyond PBM''.
\textcolor{red}{\textbf{TODO} if used to back the explicit ``PBM violations
documented in real logs'' claim: a specific quote from their experimental
chapters comparing models on real data would be needed.}

\bigskip

\section*{5. \texttt{craswell2008experimental}}
\paperclaim{Click position-bias models; the cascade model and basic position
bias hypotheses.}\\
\source{Craswell, Zoeter, Taylor, Ramsey, WSDM 2008. PDF:
\texttt{citations/craswell\_2008\_wsdm.pdf}.}

\quoted{From the abstract: ``Search engine click logs provide an invaluable
source of relevance information, but this information is biased. A key
source of bias is presentation order: the probability of click is
influenced by a document's position in the results page. [...] We
propose four simple hypotheses about how position bias might arise. [...]
\textbf{A `cascade' model, where users view results from top to bottom and
leave as soon as they see a worthwhile document, is our best explanation
for position bias in early ranks.}''}

\quoted{\S2, ``Cascade Model'' (p.~88): ``Now we present a new model for
explaining the position effect, which assumes a linear traversal through
the ranking, and that documents below a clicked result are not examined.
[...] In the cascade model, we assume that the user views search results
from top to bottom, deciding whether to click each result before moving
to the next. Each document $d_i$ is either clicked with probability $r_d$
or skipped with probability $(1-r_d)$. In the most basic form of the
model, we assume that a user who clicks never comes back, and a user who
skips always continues, in which case:
$c_{di} = r_d \prod_{j=1}^{i-1}(1 - r_{\text{doc in rank }j})$.''}

\noindent Verified. The cascade model is introduced here and matches our
usage as a PBM-violation alternative click model.
(``trust bias'' and ``position-dependent relevance'' are formalised
elsewhere -- our cite of Craswell here is purely for the cascade model
definition.)

\bigskip

\section*{6. \texttt{hesterberg1995}}
\paperclaim{Self-normalised weighted ratio estimators are standard in
importance sampling; heterogeneous-weight importance sampling variance
analysis.}\\
\source{Tim Hesterberg, ``Weighted Average Importance Sampling and
Defensive Mixture Distributions'', Technometrics 37(2):185--194, 1995.
PDF: \texttt{citations/hesterberg\_1995\_technometrics.pdf}.}

\quoted{From the abstract: ``Importance sampling uses observations from
one distribution to estimate for another distribution by weighting the
observations. Including the target distribution as one component of a
mixture distribution bounds the weights and makes importance sampling
more reliable. \textbf{The usual importance-sampling estimate is a
weighted average with weights that do not sum to 1.} We discuss simple
normalization and other, more efficient normalization methods.''}

\quoted{\S1 Introduction (p.~185): ``The integration estimate is a
\emph{weighted average with weights that do not sum to 1.} Let
$W(x) = f(x)/g(x)$ be the weight function, based on the likelihood ratio
between the target and design distributions; then the empirical
distribution function corresponding to the integration estimate has
weight $V_{int,i} = n^{-1}W_i$ on the output $Q(X_i)$ from the $i$th
replication, where $W_i = W(X_i)$. The sum of the empirical weights is
$\bar W$, which has expected value 1 but nonzero variance.''}

\noindent Verified. The paper is exactly about the variance of
self-normalised weighted IS estimators, which is what we cite for.

\bigskip

\section*{7. \texttt{joachims2017unbiased}}
\paperclaim{Position bias systematically corrupts click signals used for
learning-to-rank; IPS is the standard counterfactual estimator.}\\
\source{Joachims, Swaminathan, Schnabel, ``Unbiased Learning-to-Rank with
Biased Feedback'', WSDM 2017. PDF:
\texttt{citations/joachims\_2017\_wsdm.pdf}.}

\quoted{From the abstract: ``While implicit feedback has many advantages
(e.g., it is inexpensive to collect, user centric, and timely), \textbf{its
inherent biases are a key obstacle to its effective use. For example,
position bias in search rankings strongly influences how many clicks a
result receives, so that directly using click data as a training signal in
Learning-to-Rank (LTR) methods yields sub-optimal results.} To overcome
this bias problem, we present a counterfactual inference framework that
provides the theoretical basis for unbiased LTR via Empirical Risk
Minimization despite biased data.''}

\quoted{\S2 Related Work: ``Our approach uses inverse propensity scoring
(IPS), originally employed in survey sampling [7] and causal inference
from observational studies [19], but more recently also in whole page
optimization [29], IR evaluation with manual judgments [13, 21], and
recommender evaluation [13, 12].''}

\noindent Verified. Cite supports both ``position bias systematically
corrupts click signals'' (intro) and ``IPS / counterfactual framework
for ULTR'' (the standard reference).

\bigskip

\section*{8. \texttt{kddcup2012} (dataset, no PDF)}
\paperclaim{Dataset reference for the KDD Cup 2012 Track 2 click log.}\\
\source{Tencent / KDD Cup 2012 competition release.
\url{https://www.kaggle.com/c/kddcup2012-track2}}

\quoted{(Dataset citation; no paper. The dataset statistics in the paper
were computed by the authors directly from the released training file.)}

\bigskip

\section*{9. \texttt{kong1992ess}}
\paperclaim{Effective sample size $\mathrm{ESS}(\omega) = (\sum_i\omega_i)^2
/ \sum_i\omega_i^2$.}\\
\source{Kong, ``A note on importance sampling using standardized weights'',
Tech. Rep. 348, U.\ Chicago Statistics, 1992.
PDF: \texttt{citations/kong\_1992\_tr348.pdf}. Also re-derived in
Owen 2013 §9.3 Eq.~(9.13).}

\quoted{From Owen Ch.~9, Eq.~(9.13) (which cites Kong 1992): ``Setting
$\mathrm{Var}(S_w) = \sigma^2/n_e$ and solving for $n_e$ yields the
\textbf{effective sample size}
$n_e = (\sum_{i=1}^n w_i)^2 / \sum_{i=1}^n w_i^2 = n\bar w^2 / \overline{w^2}$.''}

\noindent Verified via Owen \S9.3.

\bigskip

\section*{10. \texttt{owen2013mc}}
\paperclaim{(i) The textbook closed-form asymptotic variance of the
self-normalised weighted ratio estimator $\widetilde V(\omega)$;
(ii) Effective sample size.}\\
\source{Owen, \emph{Monte Carlo Theory, Methods and Examples}, Ch.~9
(``Importance sampling''), \S9.2 ``Self-normalized importance sampling''
(pp.~8--10) and \S9.3 ``Importance sampling diagnostics''
(pp.~11--12). PDF: \texttt{citations/owen\_mc\_ch09\_importance.pdf}.}

\quoted{Owen \S9.2, Eq.~(9.6)--(9.9): ``Consider the \textbf{self-normalized
importance sampling estimate}
$\tilde\mu_q = \sum_i f(X_i)w_u(X_i) / \sum_i w_u(X_i)$. To place a
confidence interval around $\tilde\mu_q$ we use the delta method and
equation (2.30) from \S2.7 for the variance of a ratio of means. Applying
that formula to (9.7) yields an approximate variance for $\tilde\mu_q$ of
$\widehat{\mathrm{Var}}(\tilde\mu_q) = (1/n)\,\mathbb E_q\bigl((f(X)w(X)-\mu w(X))^2\bigr)
/ \mathbb E_q(w(X))^2 = \sigma^2_{q,sn}/n$. For a computable variance
estimate, we employ $w_u$ instead of $w$, getting
$\widehat{\mathrm{Var}}(\tilde\mu_q) = \sum_i w_i^2 (f(x_i)-\tilde\mu_q)^2$.''}

\quoted{Owen \S9.3, Eq.~(9.13): ``Setting $\mathrm{Var}(S_w) = \sigma^2/n_e$
and solving for $n_e$ yields the \textbf{effective sample size}
$n_e = (\sum_{i=1}^n w_i)^2 / \sum_{i=1}^n w_i^2$.''}

\noindent Verified.

\bigskip

\section*{11. \texttt{swaminathan2015snips}}
\paperclaim{SNIPS is the off-policy analogue of our plug-in $\widetilde V$
framework: same self-normalised weighted ratio variance applied to
counterfactual evaluation.}\\
\source{Swaminathan \& Joachims, ``The Self-Normalized Estimator for
Counterfactual Learning'', NeurIPS 2015. PDF:
\texttt{citations/swaminathan\_joachims\_2015\_nips\_snips.pdf}.}

\quoted{From the abstract: ``This paper identifies a severe problem of the
counterfactual risk estimator typically used in batch learning from
logged bandit feedback (BLBF), and \textbf{proposes the use of an
alternative estimator that avoids this problem.} [...] We show that this
conventional estimator suffers from a \emph{propensity overfitting}
problem when used for learning over complex hypothesis spaces. We
propose to replace the risk estimator with a \textbf{self-normalized
estimator}, showing that it neatly avoids this problem.''}

\quoted{\S2 Related Work (p.~2): ``Classic variance reduction techniques
for importance sampling are also useful for counterfactual evaluation and
learning. For instance, importance weights can be ``clipped'' to trade-off
bias against variance in the estimators. [...] In particular, \textbf{we
study the self-normalized estimator which is superior to the vanilla
estimator when fluctuations in the weights dominate the variance.} We
additionally show that the self-normalized estimator neatly addresses
propensity overfitting.''}

\noindent Verified. SNIPS is exactly the self-normalised IPS estimator
whose variance behaviour mirrors our $\widetilde V$ framework, and the
key insight ``superior when fluctuations in the weights dominate the
variance'' directly matches our point that weight heterogeneity drives
precision.

\bigskip

\section*{12. \texttt{vdv1998asymptotic}}
\paperclaim{Delta method.}\\
\source{van der Vaart, \emph{Asymptotic Statistics}, Cambridge 1998, Ch.~3
``Delta method''. Frontmatter PDF only:
\texttt{citations/vdv\_1998\_frontmatter.pdf}. Ch.~3 content not in our
local copy.}

\quoted{(Citation is the standard reference for the delta method; we
verify only author/title/publisher from the frontmatter. Wasserman 36-705
lec~3 \texttt{citations/wasserman\_36705\_lec03.pdf} provides a freely
available equivalent.)}

\bigskip

\section*{13. \texttt{wang2018position}}
\paperclaim{(i) PBM EM joint estimation of relevance and propensities;
(ii) PBM violations on real personal-search logs.}\\
\source{Wang, Golbandi, Bendersky, Metzler, Najork, ``Position Bias
Estimation for Unbiased Learning to Rank in Personal Search'', WSDM 2018.
PDF: \texttt{citations/wang\_2018\_wsdm.pdf}.}

\quoted{From the abstract: ``We propose a \textbf{regression-based
Expectation-Maximization (EM) algorithm} that is based on a position bias
click model and that can handle highly sparse clicks in personal search.
We evaluate our EM algorithm and the extracted bias in the
learning-to-rank setting.''}

\quoted{\S2 Related Work: ``Recently, Wang et al.\ [31] and Joachims et
al.\ [20] introduced the unbiased learning-to-rank framework by treating
the bias as a counterfactual effect. Inside the framework, the critical
component is to quantify the click bias, rather than estimate the
$\langle$query, document$\rangle$ relevance in click models.''}

\noindent (i) Verified: EM estimation of position bias is the core
contribution.

\textbf{(ii) Result of careful read of \S5--6:} the paper compares
their EM estimator's position-bias estimate to a RandPair-experiment
ground truth (\S5.2.1, Fig.~2) and shows ``Empirical has more skew than
RandPair, consistent with previous findings'' -- i.e.\ raw click rates
\emph{differ from} PBM-via-RandPair ground truth, which is evidence of
position bias but \emph{not} explicitly of PBM model mis-specification.
The Conclusions (\S6) summarise: ``the propensity estimation is
equivalent to the position bias estimation in the \emph{classical
position bias model}'' -- they work \emph{within} PBM, not against it.

\textcolor{red}{\textbf{Recommendation:} Wang 2018 is correctly cited for
EM-style joint relevance/propensity estimation but does \emph{not}
provide direct empirical evidence of PBM violations. We should restrict
the cite to the EM-comparison context in our paper and not include it in
the ``PBM violations on real logs are documented in [\,...]'' list. The
remaining cites in that list (\texttt{yue2010beyond},
\texttt{chuklin2015clickmodels}) are sufficient.}

\bigskip

\section*{14. \texttt{yue2010beyond}}
\paperclaim{Documents that real click data has effects beyond position
(attractiveness, trust bias, etc.) -- i.e.\ PBM-violating structure.}\\
\source{Yue, Patel, Roehrig, ``Beyond Position Bias: Examining Result
Attractiveness as a Source of Presentation Bias in Clickthrough Data'',
WWW 2010. PDF: \texttt{citations/yue\_2010\_www\_beyond.pdf}.}

\quoted{From the abstract: ``Although data is plentiful, one must take
care when interpreting clicks, since user behavior can be affected by
various sources of presentation bias. \textbf{While the issue of position
bias in clickthrough data has been the topic of much study, other
presentation bias effects have received comparatively little attention.}
For instance, since users must decide whether to click on a result based
on its summary (e.g., the title, URL and abstract), one might expect
clicks to favor `more attractive' results. In this paper, we examine
result summary attractiveness as a potential source of presentation
bias. [...] \textbf{Our experiments conducted on the Google web search
engine show substantial evidence of presentation bias in clicks towards
results with more attractive titles.}''}

\quoted{\S1 Introduction: ``[T]he so called position bias effect. But
users must also judge relevance based on summaries rather than the
actual results themselves. [...] Intuitively, we might expect click
behavior to favor the more attractive title. \textbf{Thus, even in the
absence of position bias, a result's perceived relevance might noticeably
differ from its actual relevance.}''}

\noindent Verified. Paper directly documents PBM-violating structure
(attractiveness biases) in real Google search logs, supporting our cite.

\bigskip

\section*{15. \texttt{zou2022large}}
\paperclaim{Baidu-ULTR is a large search dataset for unbiased
learning-to-rank; the released aggregation has $94\%$ of cells with
$N=1$ and $78\%$ of adjacent-pair ratios in $[0.5,2]$, which we
attribute to the per-cell deduplication in the public release.}\\
\source{Zou, Mao, Chu, Tang, Ye, Wang, Yin, ``A Large Scale Search
Dataset for Unbiased Learning to Rank'', NeurIPS 2022 Datasets and
Benchmarks Track. PDF: \texttt{citations/zou\_2022\_neurips.pdf}.}

\quoted{Table 1 (p.~4): ``Baidu: \# Query = 383{,}429{,}526; \# Doc =
1{,}287{,}710{,}306; \# User Feedback = 18 (Real Feedback); \# Display-info
= 8 (Display Info); \# Session = 1{,}210{,}257{,}130; Pub-Year = 2022.''}

\quoted{\S3.1 ``Documents'' (p.~4): ``For the candidate documents of the
query, we only record the displayed documents to save the cost of
storing the logs. Typically, the document not being displayed is less
informative than the ones displayed since the user cannot provide any
feedback on the document without display. As a result, the logged search
session usually contains just 10 documents for every query because one
page contains 10 results.''}

\noindent The dataset's size and the largest-public-ULTR-dataset claim
are verified. \textbf{The specific $94\%$ $N=1$ and $78\%$ ratio
statistics are not stated in Zou 2022 -- we computed those ourselves
from the released aggregation.} The cite as used in our paper supports
``Baidu-ULTR aggregation'' as a public benchmark; the deduplication
attribution is our own characterisation, not a Zou quote. This is fine
since the in-text language says ``we attribute to...''.

\end{document}
